% Part: methods
% Chapter: induction
% Section: strong-induction

\documentclass[../../../include/open-logic-section]{subfiles}

\begin{document}

\olfileid[pt]{mth}{ind}{str}

\olsection{Indução Forte}

No princípio de indução discutido acima, provamos $P(0)$ e
também que, se $P(k)$, então $P(k+1)$.  Na segunda parte, assumimos que
$P(k)$ é verdadeiro e usamos essa suposição para provar $P(k+1)$.
De forma equivalente, é claro, poderíamos assumir $P(k-1)$ e usá-la para
provar $P(k)$---o importante é que sejamos capazes de realizar a
inferência de qualquer número para seu sucessor; que possamos provar a
afirmação em questão para qualquer número sob a suposição de que ela
vale para seu antecessor.

Há uma variante do princípio de indução em que não apenas assumimos que
a afirmação vale para o antecessor $k-1$ de $k$, mas para todos os
números menores que~$k$, e usamos essa suposição para estabelecer a
afirmação para~$k$. Isso também nos dá a afirmação $P(n)$ para todo~$n
\in \Nat$.  Pois, uma vez que tenhamos estabelecido $P(0)$, teremos com
isso estabelecido que $P$ vale para todos os números menores que~$1$.
E se sabemos que, se $P(l)$ para todo $l<k$, então $P(k)$, sabemos isso
em particular para $k=1$.  Então podemos concluir $P(1)$.  Com isso
provamos $P(0)$ e $P(1)$, isto é, $P(l)$ para todo $l<2$, e como temos
também o condicional, se $P(l)$ para todo $l<2$, então $P(2)$, podemos
concluir~$P(2)$, e assim por diante.

De fato, se conseguirmos estabelecer o condicional geral ``para todo
$k$, se $P(l)$ para todo $l<k$, então $P(k)$,'' não temos mais de
estabelecer $P(0)$, já que ele decorre disso.  Pois lembre-se de que
uma afirmação geral como ``para todo $l<k$, $P(l)$'' é verdadeira se não
houver nenhum $l<k$.  Este é um caso de quantificação vazia: ``todos os
$A$s são $B$s'' é verdadeiro se não houver nenhum $A$,
$\lforall[x][(!A(x) \lif !B(x))]$ é verdadeiro se nenhum $x$
satisfaz~$!A(x)$. Neste caso, a versão formalizada seria
``$\lforall[l][(l < k \lif P(l))]$''---e isso é verdadeiro se não houver
nenhum $l < k$.  E se $k=0$ esse é exatamente o caso: não há nenhum
$l<0$, logo ``para todo $l<0$, $P(0)$'' é verdadeiro, qualquer que
seja $P$.  Uma prova de ``se $P(l)$ para todo $l<k$, então $P(k)$''
estabelece assim automaticamente~$P(0)$.

Esta variante é útil se estabelecer a afirmação para~$k$ não puder ser
feito de modo a apenas se apoiar na afirmação para $k-1$, mas puder
exigir a suposição de que ela é verdadeira para um ou mais $l<k$.

\end{document}
